\documentclass[10pt,aspectratio=169]{beamer}

\usetheme[progressbar=frametitle]{metropolis}
\usepackage{appendixnumberbeamer}
\usepackage{booktabs}
\usepackage[scale=2]{ccicons}
\usepackage{xspace}
\usepackage{amsmath}
\usepackage{amsfonts}
\usepackage{amssymb}
\usepackage{listings}
\usepackage{xcolor}

\definecolor{lightgray}{rgb}{0.95,0.95,0.95}
\definecolor{keywordcolor}{rgb}{0.1,0.1,0.8}
\definecolor{commentcolor}{rgb}{0.3,0.5,0.3}

\lstset{
    language=Python,
    basicstyle=\ttfamily\small,
    keywordstyle=\color{keywordcolor}\bfseries,
    commentstyle=\color{commentcolor}\itshape,
    backgroundcolor=\color{lightgray},
    frame=single,
    rulecolor=\color{lightgray},
    breaklines=true,
    showstringspaces=false,
    tabsize=4
}

\title{The Architecture of Truth}
\subtitle{How the Lean 4 Kernel Validates Mathematical Proofs via Dependent Type Checking}
\date{\today}
\author{A Ground-Up Reconstruction Guide}
\institute{SLP: Prof. Sudarshan Vijay}

\begin{document}

\maketitle

\begin{frame}{Agenda}
    \tableofcontents
\end{frame}

% -----------------------------------------------------------------------------
\section{1. The Type Checking Paradox}

\begin{frame}[fragile]{The Professor's Puzzle}
    \textbf{Question:} \\
    If \verb|2 + 2| is an integer, and \verb|4| is an integer... then doesn't \verb|2 + 3| also evaluate to an integer? 
    
    If they all have the same type (\verb|int|), how does "Type Checking" actually prove that \verb|2 + 2 = 4| is true, while \verb|2 + 3 = 4| is false?
\end{frame}

\begin{frame}[fragile]{The Simple Type Theory View}
    In standard programming (Python, Java, C++), type systems are strictly categorical.
    
    \begin{lstlisting}[language=Python]
print(type(2 + 2) == type(4)) # True
print(type(2 + 3) == type(4)) # True
    \end{lstlisting}
    
    A standard compiler merely validates memory bounds: "Are you passing a string to a function that requires a number?" 
    
    Under Simple Type Theory, \textbf{the professor is entirely correct}: \verb|2+3=4| is a perfectly valid, flawlessly typed expression!
\end{frame}

\begin{frame}[fragile]{Why Simple Types Fail in Mathematics}
    Mathematical theorems are not simply labels mapping variables to memory chunks. They represent exact structural relationships.
    
    If \verb|A = B| just meant "A and B have the same categorical label", the compiler would have no capability to differentiate truth from falsehood.
    
    To solve this, Lean abandons Simple Type Theory entirely.
\end{frame}

\begin{frame}[fragile]{The Solution: Dependent Type Theory}
    In Lean, the theorem \verb|2 + 2 = 4| is \textbf{not} of type \verb|int| (or \verb|Bool|).
    
    \vspace{1em}
    \textbf{The Theorem IS the Type.}
    
    \vspace{1em}
    Specifically, the theorem is a \textbf{Dependent Type} parameterized by the exact computational expressions taking place: \verb|Eq (add 2 2) (4)|.
    
    To pass type checking, the Kernel does not check the argument labels. It checks if your Proof Object successfully inhabits this explicit, mathematically unrolled Type!
\end{frame}

\begin{frame}[fragile]{Analogy: Buckets vs. Custom Puzzle Pieces}
    \textbf{Standard Programming (Simple Types):} \\
    Types are just \textbf{Buckets}. "Is this data in the Integer bucket?" Yes. \verb|2+2| goes in the bucket. \verb|2+3| goes in the bucket.
    
    \vspace{1em}
    \textbf{Lean (Dependent Types):} \\
    Types are uniquely cut \textbf{Puzzle Holes}. \\
    The Theorem \verb|2+2=4| is uniquely shaped hole. Your proof is a puzzle piece. \\
    "Type Checking" is literally just the compiler checking: \textit{"Does the specific puzzle piece you built perfectly click into this exact mathematical hole?"}
\end{frame}

% -----------------------------------------------------------------------------
\section{2. The Two-Tier Architecture}

\begin{frame}[fragile]{The de Bruijn Criterion}
    The defining architectural philosophy of Lean 4 (and Coq/Agda) is the \textbf{de Bruijn Criterion}.
    
    The system is strictly bifurcated into two drastically asymmetrical halves:
    \begin{itemize}
        \item \textbf{The Untrusted Frontend} (Massive, complex, written in Lean)
        \item \textbf{The Trusted Kernel} (Minimal, ~5,000 lines of highly optimized C++)
    \end{itemize}
\end{frame}

\begin{frame}[fragile]{The Untrusted Frontend}
    The frontend is hundreds of thousands of lines of code. It includes:
    \begin{itemize}
        \item The Parser and Elaborator
        \item Macro Expansion
        \item The Tactic Engine (\verb|rw|, \verb|induction|, etc.)
        \item Implicit Argument Inference
    \end{itemize}
    
    Crucially: The untrusted frontend is \textbf{expected} to contain bugs. It might crash or loop infinitely. It might even accidentally declare a false theorem "proven"!
\end{frame}

\begin{frame}[fragile]{Analogy: The Frontend is the Lawyer}
    Think of the Untrusted Frontend as a \textbf{lawyer}.
    
    \vspace{1em}
    The lawyer researches precedents, builds a massive argument, uses automation tools, and sometimes makes logical leaps or mistakes trying to win the case. 
    
    \vspace{1em}
    But the lawyer doesn't decide what is true. They just compile a massive, final legal brief (the Proof Term) to submit to the judge.
\end{frame}

\begin{frame}[fragile]{The Trusted Kernel}
    The Kernel is a monolithic C++ core (\verb|src/kernel|). It performs exactly one task: \textbf{Type Checking}.
    
    The frontend's only job is to generate a fully explicit, unambiguous "Proof Term" in a minimal core calculus, and hand it to the Kernel.
    
    If the Kernel type-checks it, the theorem is proven. If the frontend suffers a bug and generates a flawed proof, the Kernel ruthlessly rejects it.
\end{frame}

\begin{frame}[fragile]{Analogy: The Kernel is the Supreme Court}
    The Kernel is the \textbf{Supreme Court Justice}.
    
    \vspace{1em}
    The Justice takes the lawyer's brief. The Justice doesn't care \textit{how} the lawyer came up with the argument (automation, macros, guessing). 
    
    \vspace{1em}
    The Justice painstakingly reads the brief line-by-line, mechanically verifying if every single deductive step strictly follows the fundamental laws of math. If the lawyer made a mistake, the Justice throws the case out immediately.
\end{frame}

\begin{frame}[fragile]{The Boundary: Why Bifurcate?}
    To reconstruct Lean, you only need to build the Kernel correctly. 
    
    If the Kernel is sound, the entire system inherits that soundness, regardless of what chaos happens in IDE integrations, automation scripts, or metaprogramming.
    
    This rigorous containment model is why Lean holds absolute authority on mathematical truth.
\end{frame}

% -----------------------------------------------------------------------------
\section{3. The Core Calculus (CIC)}

\begin{frame}[fragile]{Calculus of Inductive Constructions (CIC)}
    Lean's Kernel is based on the \textbf{Calculus of Inductive Constructions}.
    
    There is no fundamental distinction between types, functions, and values. Everything is an Expression (\verb|Expr|).
\end{frame}

\begin{frame}[fragile]{Unifying Types and Values}
    In CIC:
    \begin{itemize}
        \item \verb|3| is an \verb|Expr|
        \item \verb|Nat| (the type of 3) is an \verb|Expr|
        \item \verb|Type| (the type of Nat) is an \verb|Expr|
        \item \verb|2 + 2 = 4| is an \verb|Expr| (a proposition type)
        \item A proof indicating \verb|2 + 2 = 4| is an \verb|Expr|
    \end{itemize}
\end{frame}

\begin{frame}[fragile]{The Curry-Howard Correspondence}
    This unification is the heartbeat of modern theorem proving:
    
    \vspace{1em}
    \begin{center}
        \Large \textbf{Proving a Theorem P is literally identical to constructing a program whose Type evaluates to P.}
    \end{center}
\end{frame}

\begin{frame}[fragile]{Analogy: The Mathematical Factory}
    Think of Curry-Howard as a \textbf{Factory}. 
    
    \vspace{1em}
    You want a product (a Theorem, like \verb|1 <= 4|). 
    
    Instead of writing the theorem on a whiteboard, you build a physical machine (a Program / Function) that strictly manufactures that product. 
    
    \vspace{1em}
    If the Factory Inspector (Type Checker) turns your machine on, and an object of type \verb|1 <= 4| successfully comes out the conveyor belt without jamming... you have proven the theorem!
\end{frame}

\begin{frame}[fragile]{The Hierarchy of Universes}
    If \verb|Type : Type|, the system triggers Girard's paradox (a type-theory Russell's paradox), allowing proofs of \verb|False|.
    
    Lean prevents this using an infinite hierarchy of universes (\verb|Level|):
    \begin{itemize}
        \item \verb|Sort 0| represents \verb|Prop|
        \item \verb|Sort 1| represents \verb|Type 0|
        \item \verb|Sort (u+1)| represents \verb|Type u|
    \end{itemize}
\end{frame}

\begin{frame}[fragile]{Prop vs Type}
    \begin{itemize}
        \item \textbf{\texttt{Prop} (Sort 0):} The universe of mathematical propositions (\verb|2+2=4|). These undergo "Proof Erasure"—computationally, the content of the proof doesn't matter, only that it exists.
        \item \textbf{\texttt{Type 0} (Sort 1):} The universe of computational data (\verb|Nat|, \verb|Bool|, \verb|String|). These survive compilation to execute at runtime.
    \end{itemize}
\end{frame}

\begin{frame}[fragile]{Impredicativity of Prop (\texttt{IMax})}
    \verb|Prop| is \textbf{impredicative}. 
    
    If \verb|B| is a proposition, then the function type \verb|A|\(\rightarrow\)\verb|B| is *also* a proposition, regardless of how massively huge the universe of \verb|A| is! 
    
    (e.g., "For all functions in $\mathbb{R}$, $f$ is continuous" is just a standard proposition, despite quantifying over a massive set).
    
    The Kernel's \verb|IMax| operator natively encodes this rule to securely construct functional types.
\end{frame}

% -----------------------------------------------------------------------------
\section{4. The Expr Abstract Syntax Tree}

\begin{frame}[fragile]{Lean's Internal Fabric}
    Everything in the Kernel is a node in the \verb|Expr| Abstract Syntax Tree.
    
    Lean aggressively memoizes these nodes. If two \verb|Expr| branches are structurally identical, they physically share the exact same hashed pointer in memory!
\end{frame}

\begin{frame}[fragile]{The \texttt{Expr} Data Structure}
    The Kernel operates on precisely these pure variants:
    \begin{itemize}
        \item \verb|BVar|: Bound Variable
        \item \verb|FVar|: Free Variable
        \item \verb|Const|: Global Declaration Reference
        \item \verb|App|: Function Application
        \item \verb|Lam|: Lambda Abstraction (\verb|fun x => t|)
        \item \verb|Forall|: Dependent Tye (\verb|(x:A) -> B|)
        \item \verb|Sort|: Universe Level
    \end{itemize}
\end{frame}

\begin{frame}[fragile]{The Variable Problem ($\alpha$-equivalence)}
    A massive system challenge: \verb|fun x => x| is mathematically identical to \verb|fun y => y|. But string names make them look different.
    
    Accidental variable capture during substitution is the \#1 cause of bugs in naive theorem provers.
    
    Lean solves this by abandoning string names entirely.
\end{frame}

\begin{frame}[fragile]{Bound Variables \& De Bruijn Indices}
    \verb|BVar| nodes use \textbf{De Bruijn Indices}. They contain no name. They merely contain an integer pointer tracking how many lambda-binders to "jump over" going upward in the AST.
    
    \begin{itemize}
        \item \verb|fun x => fun y => x (y x)|
        \item Becomes: \verb|lam( lam( App( BVar(1), App(BVar(0), BVar(1)) ) ) )|
    \end{itemize}
    
    $\alpha$-equivalence instantly becomes identical to raw pointer memory hashing!
\end{frame}

\begin{frame}[fragile]{Analogy: GPS Directions instead of Names}
    Imagine telling the computer: "Go to John's house." The computer says: "Who is John? Are there two Johns?"
    
    \vspace{1em}
    De Bruijn indices stop using names and use \textbf{GPS Directions}.
    
    \vspace{1em}
    "Go to the variable defined exactly $N$ layers above me." 
    
    Because the route to a variable is physically encoded, two identical equations will always map to the exact same GPS coordinates, proving they are the same equation!
\end{frame}

\begin{frame}[fragile]{Free Variables (\texttt{FVar})}
    \verb|FVar|s represent variables in the local proof context (hypotheses) that are \textbf{not} bound by an enclosing lambda.
    
    They use globally unique generated UUIDs. This completely mathematically nullifies any chance of accidental capture during substitution by the Kernel.
\end{frame}

\begin{frame}[fragile]{Metavariables (\texttt{MVar})}
    \verb|MVar|s (\verb|?m|) represent "holes" that the Frontend Elaborator is trying to guess (unification).
    
    \vspace{1em}
    \textbf{Crucial Security Rule:} \\
    The Trusted Kernel explicitly \textbf{bans} Metavariables! Before an expression is submitted for type checking, the frontend must instantiate all holes.
\end{frame}

% -----------------------------------------------------------------------------
\section{5. Definitional Equality \& Reductions}

\begin{frame}[fragile]{The Core Typechecker API (\texttt{inferType})}
    The absolute core execution loop of the Kernel. Given an Environment and an \verb|Expr|, it computes the type.
    
    \vspace{1em}
    Signature: \verb|Expr inferType(Expr e, LocalContext lctx)|
    
    When evaluating an \verb|App(f, a)|, it computes the type of \verb|f|, the type of \verb|a|, and checks if the argument precisely matches the parameters requirement.
\end{frame}

\begin{frame}[fragile]{Syntactic vs Definitional Equality}
    In standard programming, \verb|Array(2+2)| and \verb|Array(4)| trigger type errors because \verb|"2+2"| != \verb|"4"| syntactically.
    
    Lean utilizes \textbf{Definitional Equality} (\verb|isDefEq|). It recognizes that computed types can contain logic. It must forcefully compute both sides into a normal form to see if they match structurally!
\end{frame}

\begin{frame}[fragile]{Analogy: Grinding Coffee Beans}
    Imagine comparing a bag of "Whole Roast Beans" and a cup of "Liquid Coffee". They look completely different syntactically.
    
    \vspace{1em}
    But Lean's \verb|isDefEq| engine puts both of them through the computational \textbf{grinder}. 
    
    \vspace{1em}
    It aggressively grinds them down to their absolute basic molecular parts. Once they are ground down into primitive molecules, the Kernel puts them under a microscope to see if the cells match mathematically.
\end{frame}

\begin{frame}[fragile]{The \texttt{isDefEq} Engine}
    To determine if \verb|Type A == Type B|, the Kernel possesses a Turing-complete Evaluation Engine.
    
    It incrementally unleashes four distinct Reduction routines in parallel until a structural hash-match is found.
\end{frame}

\begin{frame}[fragile]{The Four Reductions: $\beta$ (Beta) \& $\delta$ (Delta)}
    \begin{itemize}
        \item \textbf{$\beta$-Reduction:} Function application execution. 
        \\ \verb|App( Lam(x, t), a )| $\implies$ \verb|t[x := a]|
        
        \vspace{1em}
        \item \textbf{$\delta$-Reduction:} Unfolding transparent definitions.
        \\ \verb|double 2| $\implies$ \verb|2 + 2|
    \end{itemize}
\end{frame}

\begin{frame}[fragile]{The Four Reductions: $\zeta$ (Zeta) \& $\iota$ (Iota)}
    \begin{itemize}
        \item \textbf{$\zeta$-Reduction:} Expanding local \verb|Let| bindings securely inline.
        
        \vspace{1em}
        \item \textbf{$\iota$-Reduction:} Evaluating structural recursion primitives over inductive data types (e.g., Peano numbers, lists).
        \\ \verb|Nat.rec f_z f_s (succ n)| $\implies$ \verb|f_s n (Nat.rec f_z f_s n)|
    \end{itemize}
\end{frame}

% -----------------------------------------------------------------------------
\section{6. Inductive Data Types}

\begin{frame}[fragile]{Declaring Data Types}
    When a developer writes \verb|inductive Nat | \verb|| zero | \verb|| succ : Nat -> Nat|, the Kernel validates it for strict positivity to prevent paradoxes.
\end{frame}

\begin{frame}[fragile]{The Powerful Recursor (\texttt{.rec})}
    Lean does not have primitive \verb|if/else| or \verb|match| statements in the Kernel!
    
    Every inductive type auto-generates a \textbf{Recursor} (\verb|Nat.rec|). Any code utilizing pattern-matching is compiled by the frontend into a massive tree of applications using this pure recursor, which the $\iota$-Reduction engine evaluates natively.
\end{frame}

\begin{frame}[fragile]{Quotients and Equivalences}
    To do calculus or modular arithmetic, we must treat mathematically distinct objects identically (e.g., Fraction $1/2 == 2/4$).
    
    Lean explicitly hardcodes \verb|Quot| systems into the Kernel. You can only project data out of an equivalence class if you provide a rigorous mathematical proof that your function respects the relation!
\end{frame}

% -----------------------------------------------------------------------------
\section{7. Resolving the Professor's Puzzle}

\begin{frame}[fragile]{Step 1: The Equation is a Type}
    Let's return to the puzzle: How does type checking prove \verb|2 + 2 = 4|, but reject \verb|2 + 3 = 4|?
    
    \vspace{1em}
    Remember, the Kernel does not check if they evaluate to \verb|int|. 
    
    It checks if the generated proof cleanly inhabits the Dependent Type:
    \begin{center}
        \Large \texttt{Eq( Add(2, 2) , 4 )}
    \end{center}
\end{frame}

\begin{frame}[fragile]{Step 2: Evaluating \texttt{Eq(Add(2,2), 4)}}
    The Kernel calls \verb|isDefEq( Add(2,2) , 4 )|.
    
    1. It unleashes $\delta$ and $\iota$ reductions on the Left Side. The Peano axioms unroll \verb|Add(2,2)| down to the primitive data tree: \verb|S(S(S(S(Z))))|.
    
    2. It subjects the Right Side to the same grinder. \verb|4| breaks down to \verb|S(S(S(S(Z))))|.
\end{frame}

\begin{frame}[fragile]{Step 3: Success! Identity Verified}
    3. The Kernel hashes the two primitive memory trees.
    
    \vspace{1em}
    \begin{center}
        \Large \texttt{S(S(S(S(Z)))) == S(S(S(S(Z))))}
    \end{center}
    
    \vspace{1em}
    They match perfectly! The Definitional Equality check passes. The type is legally inhabited. The theorem is true.
\end{frame}

\begin{frame}[fragile]{Step 4: Evaluating \texttt{Eq(Add(2,3), 4)}}
    Now, what about the false theorem? The Kernel calls \verb|isDefEq( Add(2,3) , 4 )|.
    
    1. It unleashes reductions on the Left Side. \verb|Add(2,3)| ruthlessly grinds down into \verb|S(S(S(S(S(Z)))))|.
    
    2. It subjects the Right Side to the grinder. \verb|4| breaks down into \verb|S(S(S(S(Z))))|.
\end{frame}

\begin{frame}[fragile]{Step 5: Failure! The Trees Diverge}
    3. The Kernel hashes the two resulting memory trees.
    
    \vspace{1em}
    \begin{center}
        \Large \color{red} \texttt{S(S(S(S(S(Z))))) != S(S(S(S(Z))))}
    \end{center}
    
    \vspace{1em}
    The structural shapes are completely different! 
\end{frame}

\begin{frame}[fragile]{The Verdict: The Absolute Authority of the Kernel}
    Because they represent divergent physical structures in memory, \verb|isDefEq| fails.
    
    The Kernel mathematically concludes that the Type \verb|Eq(Add(2,3), 4)| is completely uninhabitable. It rejects the payload, throwing a \textbf{Type Error}.
    
    This exact mechanism—relying on pure structural reduction to primitive data bits rather than loose categorical tags—is why the Lean Kernel holds the ultimate authority on mathematical truth.
\end{frame}

\begin{frame}[standout]
    Questions?
\end{frame}

\end{document}
